\documentclass[11pt]{article}
\usepackage{amsmath,amssymb,amsthm}
\usepackage{booktabs}
\usepackage{tikz}
\usepackage{lipsum}
\usepackage{hyperref}
\newtheorem{theorem}{Theorem}
\title{Notes on Incremental Typesetting}
\author{Training Document}
\begin{document}
\maketitle
\tableofcontents

\section{Introduction}\label{sec:intro}
A document is typeset from the top every time, which is why a single changed word
near the end of a long text used to cost as much as the first run did.
Keeping the engine resident lets the next run start from the last state that the
edit could not have touched, so most keystrokes only redo a page or two.
Section~\ref{sec:model} describes the model and Section~\ref{sec:results} the
numbers.\footnote{All timings are wall clock on a laptop.}

\lipsum[1-6]

\section{Model}\label{sec:model}
Every page boundary is a checkpoint of the complete engine state, compared by
fingerprint rather than by value, so storing one costs little more than the
memory it changed since the previous checkpoint.
\begin{theorem}\label{thm:sound}
If the input before position $p$ is unchanged, the state at the last checkpoint
before $p$ equals the state a full run would reach there.
\end{theorem}
\begin{proof}
TeX reads its input strictly in order, so nothing after $p$ can influence the
state before it, except through files written by an earlier run.
\end{proof}
The cost of an edit is then
\begin{equation}\label{eq:cost}
  C(p) = \sum_{i = k(p)}^{n} t_i \leq \frac{n - k(p) + 1}{n} \, T,
\end{equation}
where $k(p)$ is the first page at or after $p$ and $T$ the time of a full run.

\lipsum[7-12]

\section{Results}\label{sec:results}
Table~\ref{tab:times} lists the time per keystroke for three documents of
different length, measured with the preview open on the edited page.
\begin{table}[ht]
  \centering
  \begin{tabular}{lrrr}
    \toprule
    Document & Pages & Full run (ms) & Keystroke (ms) \\
    \midrule
    Letter & 1 & 180 & 9 \\
    Paper & 12 & 640 & 14 \\
    Thesis & 140 & 5200 & 21 \\
    \bottomrule
  \end{tabular}
  \caption{Time per keystroke against document length.}\label{tab:times}
\end{table}
Figure~\ref{fig:growth} shows that the keystroke cost grows much more slowly than
the document, as Equation~\eqref{eq:cost} and Theorem~\ref{thm:sound} predict.
\begin{figure}[ht]
  \centering
  \begin{tikzpicture}[scale=0.8]
    \draw[->] (0,0) -- (6,0) node[right] {pages};
    \draw[->] (0,0) -- (0,4) node[above] {ms};
    \draw[thick, blue] (0,0.4) -- (1,0.6) -- (3,0.9) -- (5.5,1.2);
    \draw[thick, red, dashed] (0,0.4) -- (5.5,3.6);
  \end{tikzpicture}
  \caption{Keystroke cost (solid) against a full run (dashed).}\label{fig:growth}
\end{figure}

\lipsum[13-20]

\section{Lists and Quotations}
The remaining work falls into a few groups, roughly in order of how often a
writer runs into them during an ordinary editing session.
\begin{itemize}
  \item Edits inside a float move the float, and with it every page after it.
  \item Changes to the preamble invalidate every checkpoint at once.
  \item Bibliography and index runs depend on files written by the previous run.
\end{itemize}
\begin{quote}
A preview that lags behind the cursor is read less than one that keeps up,
however beautiful its pages are once they finally arrive.
\end{quote}
\begin{enumerate}
  \item Measure where the time goes before changing anything.
  \item Keep the output identical to a full run, page for page.
\end{enumerate}

\lipsum[21-30]

\section{Conclusion}
Resident typesetting turns the cost of an edit from the length of the document
into the distance between the edit and the end of the visible page, which for
most writing is a page or two, and that is what makes a live preview possible.

\begin{thebibliography}{9}
\bibitem{knuth} D.~E. Knuth, \emph{The \TeX book}, Addison-Wesley, 1984.
\bibitem{lamport} L.~Lamport, \emph{\LaTeX: A Document Preparation System}, 1994.
\end{thebibliography}
\end{document}
